\documentclass[10pt,twocolumn]{article}

\usepackage[a4paper,margin=0.65in]{geometry}
\usepackage{mathptmx}
\usepackage[T1]{fontenc}
\usepackage[utf8]{inputenc}
\usepackage{amsmath,amssymb}
\usepackage{enumitem}
\usepackage{setspace}
\usepackage[hidelinks]{hyperref}

\setlength{\columnsep}{0.22in}
\setlength{\parindent}{9pt}
\setlength{\parskip}{0pt}

\title{\textbf{\Large \textsc{AstraVul}: AST-Guided Retrieval-Augmented Generation for Automated Software Vulnerability Detection and Triage: A Literature Review}}

\author{
    \textbf{Abdullah Al Noman} \quad \textbf{Kazi Neyamul Hasan} \quad \textbf{Rakibul Hassan}\\
    \textbf{Mahathir Mohammad} \quad \textbf{Md. Habibulla Misba}\\[4pt]
    Department of Computer Science and Engineering\\
    United International University (UIU)
}
\date{\today}

\begin{document}

\maketitle

\begin{abstract}
\textbf{\textit{Abstract}}---Automated vulnerability detection is vital for modern software security. Traditional Static Application Security Testing (SAST) suffers from high False Discovery Rates ($\text{FDR} \ge 40\%$), while naive querying of Large Language Models (LLMs) on raw source code introduces intellectual property risks, severe context saturation, and hallucinated reports. Recent research converges on hybrid neuro-symbolic paradigms: deterministic program slicing via Abstract Syntax Trees (AST) combined with knowledge-level Retrieval-Augmented Generation (RAG) and local, privacy-preserving LLM verification. This paper synthesizes 15 foundational and state-of-the-art Q1 papers (including USENIX Security, NDSS, ICSE, ACM TOSEM, IEEE TDSC, NeurIPS, and ACL) into a unified taxonomy, identifying critical research gaps and establishing the foundation for \textbf{\textsc{AstraVul}}, a token-efficient, high-precision triage pipeline.
\end{abstract}

\vspace{2pt}
\noindent\textbf{\textit{Keywords}}---Vulnerability Detection, Abstract Syntax Tree, Program Slicing, Retrieval-Augmented Generation, Local LLM, ChromaDB, Neuro-Symbolic Security, AstraVul.

\section{Introduction}
Modern software vulnerability detection has become indispensable in the Secure Software Development Life Cycle (SSDLC). However, current solutions exhibit critical limitations:
\begin{enumerate}[leftmargin=10pt,topsep=1pt,itemsep=0.5pt]
    \item \textbf{Heuristic Static Analysis (SAST):} Pattern-matching tools (e.g., Flawfinder, Cppcheck) generate excessive false positives ($\text{FDR} \ge 40\%$) due to a lack of semantic context and dataflow awareness.
    \item \textbf{Graph-Based Deep Learning:} Graph Neural Networks (e.g., Devign~\cite{zhou2019devign}, SySeVR~\cite{li2021sysevr}) capture structural relationships but operate as uninterpretable black boxes and struggle with out-of-distribution code.
    \item \textbf{Naive Large Language Models (LLMs):} While LLMs offer powerful reasoning, recent empirical audits (e.g., PrimeVul~\cite{zhang2025howfar}, Chasing Shadows~\cite{arp2026chasingshadows}) show they suffer from ``patch-blindness'', attention degradation on large files, and enterprise data privacy risks.
\end{enumerate}
To address these challenges, we review 15 seminal Q1 papers and establish the architectural motivation for \textsc{AstraVul}, an AST-guided, RAG-grounded neuro-symbolic triage engine.

\section{Thematic Synthesis of Reviewed Literature}
We categorize the reviewed literature into three complementary technical pillars:

\vspace{2pt}\noindent\textbf{1. Program Slicing and Graph Representations:}
Early deep learning on source code was hampered by excessive syntactical noise in whole-function representations. Li et al.~\cite{li2018vuldeepecker} introduced \textit{code gadgets} in VulDeePecker, demonstrating that slicing backward dataflow around library/API call sinks significantly outperforms full-function classification. SySeVR~\cite{li2021sysevr} generalized this by defining Syntax-, Semantics-, and Vulnerability-related Slices (SeVCs) across Control Flow Graphs (CFG) and Data Flow Graphs (DFG), detecting 15 zero-day vulnerabilities in open-source software. Devign~\cite{zhou2019devign} established that ASTs alone are insufficient; composite Code Property Graphs (CPGs) combining AST, CFG, and DFG embedded via Gated Graph Neural Networks (GGNN) are necessary for complex memory flaws. DeepDFA~\cite{steenhoek2024deepdfa} enriched graph embeddings with compiler-level dataflow bit-vectors, making models resilient to semantic-preserving renaming. Recently, Lekssays et al.~\cite{lekssays2025llmxcpg} (LLMxCPG) showed that deterministic CPG slicing reduces LLM input tokens by 67.8\%--90.9\% while boosting F1-scores by 15\%--40\%.

\vspace{2pt}\noindent\textbf{2. Benchmark Realities and Empirical LLM Fragility:}
Initial high accuracy claims in vulnerability detection were severely undermined by benchmark artifacts. ReVeal~\cite{chakraborty2021reveal} revealed that models trained on synthetic benchmarks (e.g., NIST Juliet) overfit to artificial syntax, experiencing catastrophic precision drops (collapsing to 10\%--14\%) in production repositories. DiverseVul~\cite{chen2023diversevul} audited historical corpora (e.g., Big-Vul) and discovered up to 48\% label noise, establishing a cleaned benchmark of 18,945 real-world C/C++ vulnerabilities. CVEfixes~\cite{bhandari2021cvefixes} automated the harvesting of ground-truth fix commits from NVD, proving that contrastive pre-patch/post-patch pairs are critical. In the LLM domain, PrimeVul~\cite{zhang2025howfar} showed that open-source 7B models achieving 68.3\% F1 on legacy datasets collapsed to 3.1\% F1 on debiased real-world data. In Chasing Shadows, Arp et al.~\cite{arp2026chasingshadows} surveyed 72 security papers, formalizing 9 systemic pitfalls and demonstrating that models must be evaluated using the False Discovery Rate ($\text{FDR} = \text{FP}/(\text{TP} + \text{FP}$).

\vspace{2pt}\noindent\textbf{3. Pre-trained Embeddings, RAG, and Neuro-Symbolic Verification:}
Pre-trained transformers advanced code comprehension: UniXcoder~\cite{guo2022unixcoder} unified comments, flattened AST prefixes, and code tokens in a 768-dimensional multimodal space. LineVul~\cite{fu2022linevul} utilized CodeBERT self-attention weights to localize flaws to exact lines, reducing developer effort by $>80\%$. PDBERT~\cite{liu2024pdbert} pre-trained on Control and Data Dependency Prediction, achieving $23\times$ higher throughput than heavy static analyzers. In neuro-symbolic triage, GPTScan~\cite{sun2024gptscan} combined static candidate extraction with LLM constraint validation ($>90\%$ precision). Addressing LLM ``patch-blindness'', Vul-RAG~\cite{du2024vulrag} injected external CWE rules, CVE descriptions, and patch diffs into the prompt, boosting accuracy by 16\%--24\%.

Table~1 provides a structured comparative taxonomy synthesizing these 15 seminal works across analysis paradigm, program representation, retrieval mechanism, privacy model, and addressed limitation.

\begin{table*}[!t]
\centering
\caption{Systematic Comparative Taxonomy of Reviewed Literature}
\vspace{3pt}
\scriptsize
\setlength{\tabcolsep}{3.5pt}
\begin{tabular}{|p{2.5cm}|p{3.2cm}|p{2.0cm}|c|p{2.1cm}|p{5.6cm}|}
\hline
\textbf{Paper / System} & \textbf{Approach Paradigm} & \textbf{Slicing/AST} & \textbf{RAG} & \textbf{Privacy / Local} & \textbf{Key Addressed Limitation} \\
\hline
VulDeePecker~\cite{li2018vuldeepecker} & Deep Learning (BLSTM) & API Sinks & No & Offline Model & Coarse file-level analysis \\
SySeVR~\cite{li2021sysevr} & Deep Learning (BBiRNN) & AST+CFG+DFG & No & Offline Model & Missing control-flow context \\
Devign~\cite{zhou2019devign} & Graph Neural Net (GGNN) & Composite CPG & No & Offline Model & Syntax-only representation \\
ReVeal~\cite{chakraborty2021reveal} & Graph DL / Empirical & Graph Slices & No & Offline Model & Exposed synthetic dataset bias \\
UniXcoder~\cite{guo2022unixcoder} & Cross-Modal Pre-training & AST Flattening & No & Local Weights & Semantic code embedding gap \\
LineVul~\cite{fu2022linevul} & Transformer (CodeBERT) & Token Sequence & No & Local Weights & Coarse function-level alerts \\
CVEfixes~\cite{bhandari2021cvefixes} & Curated Security Corpus & Function/Diff & No & N/A & Ground-truth CVE-to-fix mapping \\
DiverseVul~\cite{chen2023diversevul} & Benchmark Evaluation & Function Level & No & N/A & Mitigated label noise across 150 CWEs \\
GPTScan~\cite{sun2024gptscan} & Neuro-Symbolic Hybrid & Static Analysis & No & Cloud API (GPT-3.5) & $>90\%$ precision in logic flaw verification \\
DeepDFA~\cite{steenhoek2024deepdfa} & Neuro-Symbolic / GNN & Dataflow Vectors & No & Local Weights & Injected symbolic dataflow into embeddings \\
PDBERT~\cite{liu2024pdbert} & Pre-trained Transformer & Program Dep. & No & Local Weights & $23\times$ faster dependency-aware throughput \\
Vul-RAG~\cite{du2024vulrag} & Retrieval-Augmented LLM & No & Yes & Cloud API (GPT-4) & Overcame LLM patch-blindness (+16--24\%) \\
LLMxCPG~\cite{lekssays2025llmxcpg} & Neuro-Symbolic Hybrid & CPG Slicing & No & Local Fine-Tuning & 68--91\% code reduction; +15--40\% F1 score \\
PrimeVul~\cite{zhang2025howfar} & Empirical Benchmarking & Function Level & No & Local \& Cloud LLMs & Proved raw LLMs fail on debiased datasets \\
Chasing Shadows~\cite{arp2026chasingshadows} & Systematic Review & N/A & No & N/A & Formalized 9 pitfalls; highlighted FDR metric \\
\hline
\textbf{\textsc{AstraVul} (Ours)} & \textbf{AST-Guided RAG+LLM} & \textbf{AST+DFG ($\le 150$t)} & \textbf{Yes} & \textbf{Local 4-bit Quant.} & \textbf{Resolves FDR, hallucination, \& privacy} \\
\hline
\end{tabular}
\end{table*}

\vspace{2pt}\noindent\textbf{Research Gaps and \textsc{AstraVul} Architecture:}
This synthesis exposes three central gaps: (1)~\textit{Context Saturation:} Raw file prompts dilute attention; \textsc{AstraVul}'s deterministic Tree-sitter AST slicing prunes $\ge 70\%$ of unneeded tokens into minimal $\le 150$-token slices. (2)~\textit{Patch-Blindness:} General LLMs hallucinate false alarms on patched code; grounding slices via ChromaDB indexed with NIST CWE specifications and DiverseVul/CVEfixes diffs provides explicit boundary invariants. (3)~\textit{Privacy and Hallucination:} Cloud LLM reliance risks code leaks; \textsc{AstraVul}'s local 4-bit quantized models (Qwen2.5-Coder-7B) coupled with Pydantic JSON schema decoding guarantee verifiable, zero-leakage triage reports.

\section{Conclusion}
This literature review establishes the foundation for AST-guided RAG in software vulnerability detection. Uniting deterministic AST program slicing, semantic knowledge retrieval, and local constrained LLM triage directly resolves the core trade-offs between precision, speed, explainability, and code privacy.

\small
\begin{thebibliography}{15}
\setlength{\itemsep}{0.5pt}
\setlength{\parskip}{0pt}

\bibitem{li2018vuldeepecker}
Z.~Li et al., ``VulDeePecker: A deep learning-based system for vulnerability detection,'' in \emph{Proc. NDSS}, 2018.

\bibitem{li2021sysevr}
Z.~Li et al., ``SySeVR: A framework for using deep learning to detect software vulnerabilities,'' \emph{IEEE TDSC}, 19(4):2244--2258, 2021.

\bibitem{zhou2019devign}
Y.~Zhou et al., ``Devign: Effective vulnerability identification via graph neural networks,'' in \emph{Proc. NeurIPS}, 2019.

\bibitem{chakraborty2021reveal}
S.~Chakraborty et al., ``Deep learning for vulnerability detection: Are we there yet?,'' in \emph{Proc. ICSE}, pp. 883--895, 2021.

\bibitem{fu2022linevul}
M.~Fu and C.~Tantithamthavorn, ``LineVul: A transformer-based line-level vulnerability prediction,'' in \emph{Proc. MSR}, pp. 608--620, 2022.

\bibitem{guo2022unixcoder}
D.~Guo et al., ``UniXcoder: Unified cross-modal pre-training for code representation,'' in \emph{Proc. ACL}, pp. 7212--7225, 2022.

\bibitem{bhandari2021cvefixes}
G.~Bhandari et al., ``CVEfixes: Automated collection of vulnerabilities and their fixes,'' in \emph{Proc. MSR}, pp. 241--245, 2021.

\bibitem{chen2023diversevul}
Y.~Chen et al., ``DiverseVul: A new vulnerable source code dataset for deep learning,'' in \emph{Proc. RAID}, pp. 654--668, 2023.

\bibitem{sun2024gptscan}
Y.~Sun et al., ``GPTScan: Detecting logic vulnerabilities by combining GPT with program analysis,'' in \emph{Proc. ICSE}, 2024.

\bibitem{steenhoek2024deepdfa}
B.~Steenhoek et al., ``Dataflow analysis-inspired deep learning for vulnerability detection,'' in \emph{Proc. ICSE}, 2024.

\bibitem{liu2024pdbert}
Z.~Liu et al., ``Pre-training by predicting program dependencies for vulnerability analysis,'' in \emph{Proc. ICSE}, 2024.

\bibitem{du2024vulrag}
X.~Du et al., ``Vul-RAG: Enhancing LLM-based vulnerability detection via knowledge-level RAG,'' \emph{ACM TOSEM}, 2024.

\bibitem{lekssays2025llmxcpg}
A.~Lekssays et al., ``LLMxCPG: Context-aware vulnerability detection through CPG-guided LLMs,'' in \emph{Proc. USENIX Security}, 2025.

\bibitem{zhang2025howfar}
Y.~Zhang et al., ``Vulnerability detection with code language models: How far are we?,'' in \emph{Proc. ICSE}, 2025.

\bibitem{arp2026chasingshadows}
D.~Arp et al., ``Chasing shadows: Pitfalls in LLM security research,'' in \emph{Proc. NDSS}, 2026.

\end{thebibliography}

\end{document}
